\scanpage{60}

% B1 | text | upper lines, definition of the functor category
\unsure{Alors} pour tout objet $Y$ de $\unsure{X|C}$, \deleted{$\unsure{i(X)(Y)}$} posant $F=i(X)$,
$F(Y)$ est un module \inserted{à droite} projectif de type fini sur
$o'^{\unsure{o|0}}$, où $o'=\operatorname{End}(Y)$. Soit \deleted{\unsure{Cat}}, dans le cas général,
\deleted{$\unsure{\mathcal F}$} $C^{\unsure{\circ|i}}_{\unsure{\mathrm{sfp}|\mathrm{sp}}}(\Lambda)$ la catégorie de ces foncteurs \inserted{(dits ``spéciaux'')}. On trouve donc

% B2 | margin note | crossed-out left marginal note opposite the definition
% slanted
\margin{\deleted{\unsure{Ce sont les cycles|Ce sont les objets} \unsure{de type} ``fini'' au sens \unsure{général}}}

% B3 | text | line above arrow
ici une inclusion (sous réserve d'existence des quotients)

% B4 | diagram | centred inclusion
\diagram{\begin{tikzcd}[ampersand replacement=\&]
C\arrow[r,"i"]\& C_{\mathrm{sfp}}(\Lambda)
\end{tikzcd}},

% B5 | text | paragraph below inclusion
et le choix d'un \deleted{générateur} \inserted{objet} $X$ de $C$ permet d'identifier $C_{\mathrm{sp}}(\unsure{\Lambda|M})$ à la catégorie des $o$-modules
à droite projectifs de type fini. \deleted{\unsure{On}} On l'appelle
la \deleted{\unsure{enveloppe linéaire}} \deleted{\unsure{saturé projectif}} \inserted{\unsure{simplifié|simplétion|simplétée}} (ou le saturé) de $C$.

% B6 | margin note | inclined note at lower left, alongside projectifs de type fini
% slanted
\margin{il peut être préférable de dire partout : loc. libre de rang fini sur $\operatorname{Spec}\Lambda=S$}
% archive sheet number: 59
